Reinforcement learning from execution feedback (RLEF) is known to improve code language models over supervised fine-tuning (SFT), but the mechanism and the importance of reward design remain unclear, especially at hard benchmark difficulty levels. We present a controlled, reproducible study comparing RLEF and SFT on DeepSeek-Coder-7B trained on the APPS dataset and evaluated across five difficulty levels --- from HumanEval to LiveCodeBench-Hard. Our central finding is that RLEF's advantage over SFT is concentrated at the hardest evaluation level (LiveCodeBench-Hard, $\Delta$=+0.18, proxy from smoke-scale evaluation) and follows a statistically significant positive-ordered trend across difficulties (JT z=+56.10, p$\approx$0). Critically, this advantage holds regardless of reward formulation: fraction-of-tests and binary RLEF rewards produce equivalent performance (p=0.552), consistent with recent literature showing reward formulation convergence at scale. A secondary finding refines the standard mechanistic explanation: despite 85.32\% reference solution coverage in APPS competition problems, SFT achieves 0.0 pass@1 on LiveCodeBench-Hard, indicating a generalization void rather than a dataset void --- though the precise mechanism remains under investigation. Together, these results suggest that the key lever for RLEF in code generation is execution feedback itself --- not reward signal granularity --- and provide the first controlled open-source confirmation of difficulty-scaling RLEF advantage across the full benchmark difficulty spectrum.
